\chapter{Présentation du projet}
\section{Cadre du projet}
Le projet s’inscrit dans le cadre du module \textbf{Big Data} du programme \textbf{IOT3} à la \textbf{Faculté des Sciences de Tunis (FST)}. 
Il mobilise des connaissances théoriques en finance, gestion des données et technologies Big Data, ainsi que des compétences pratiques en 
modélisation UML et visualisation des données. 

\section{Problématique}
La fraude bancaire constitue un risque majeur pour les institutions financières. 
Les méthodes traditionnelles, basées sur des règles statiques et des traitements en batch, ne permettent pas de détecter efficacement les fraudes complexes, surtout avec l’augmentation du volume de transactions 
et la diversité des comportements clients~\cite{kpmg2025}.

Le projet vise à répondre aux questions suivantes :
\begin{itemize}
    \item Comment détecter la fraude bancaire de manière efficace avant l’apparition du Big Data ?
    \item Quelles sont les limites des méthodes traditionnelles en termes de précision, rapidité et scalabilité ?
    \item Comment les technologies Big Data peuvent-elles améliorer la détection des fraudes grâce à l’analyse en temps réel et à l’apprentissage automatique ?
\end{itemize}

\section{Objectifs du projet}
Les principaux objectifs sont :
\begin{enumerate}
    \item Étudier et comparer les méthodes de détection de fraude bancaire \textbf{avant et après l’apparition du Big Data}.
    \item Concevoir un système de détection de fraude basé sur des données structurées et semi-structurées.
    \item Modéliser le système à l’aide de \textbf{diagrammes UML} (classes, cas d’utilisation, séquences).
    \item Réaliser une visualisation interactive des résultats avec des outils de \textbf{Business Intelligence}.
    \item Fournir un rapport structuré et une présentation synthétique pour valoriser les travaux réalisés.
\end{enumerate}

\section{Choix méthodologiques}

\subsection{Formalisme de conception}
Le projet utilise le \textbf{UML (Unified Modeling Language)} pour représenter :
\begin{itemize}
    \item La structure du système via le \textbf{diagramme de classes}.
    \item Les interactions entre les acteurs et le système via les \textbf{diagrammes de cas d’utilisation}.
    \item La séquence des opérations pour chaque fonctionnalité clé via les \textbf{diagrammes de séquences}.
\end{itemize}

\subsection{Démarche suivie}
La démarche adoptée est inspirée du \textbf{Scrum}, une méthode agile permettant :
\begin{itemize}
    \item Une planification progressive et flexible des tâches.
    \item Des réunions régulières pour suivre l’avancement.
    \item Une amélioration continue des livrables.
\end{itemize}

Le projet a été découpé en \textbf{sprints}, chacun correspondant à une étape clé : 
analyse du domaine, collecte et traitement des données, conception UML, développement et visualisation des résultats.

\section{Conclusion}
Ce chapitre a présenté le cadre général du projet, la problématique, les objectifs et les choix méthodologiques. 
Il prépare le lecteur à comprendre les chapitres suivants, qui détailleront l’étude comparative des méthodes de détection de fraude 
avant et après l’apparition du Big Data, la conception détaillée du système et les visualisations réalisées pour analyser les résultats.